\documentclass[tikz,border=6pt]{standalone}
\usepackage{ctex}
\usepackage{amsmath}
\usepackage{pgfplots}
\pgfplotsset{compat=1.18}
\usetikzlibrary{calc,angles,quotes,arrows.meta,positioning,decorations.markings}
\begin{document}
\begin{tikzpicture}[scale=3.2,>=Latex]
  % axes
  \draw[->,gray!70] (-1.35,0) -- (1.45,0) node[right,black] {$\mathrm{Re}$};
  \draw[->,gray!70] (0,-1.4) -- (0,1.45) node[above,black] {$\mathrm{Im}$};
  % unit circle
  \draw[thick,blue!55] (0,0) circle (1);
  \coordinate (O) at (0,0);
  % five fifth-roots of unity at 2*pi*k/5  -> degrees 72*k
  \foreach \k in {0,1,2,3,4} {
    \coordinate (z\k) at ({cos(72*\k)},{sin(72*\k)});
  }
  % regular pentagon
  \draw[very thick,red!70,fill=red!8] (z0) -- (z1) -- (z2) -- (z3) -- (z4) -- cycle;
  % radii (dashed) to show equal angular spacing
  \foreach \k in {0,1,2,3,4} {
    \draw[dashed,gray!60] (O) -- (z\k);
  }
  % angle arc for spacing 2pi/5 between z0 and z1
  \draw[->,thick,purple] (0.32,0) arc (0:72:0.32);
  \node[purple] at ({0.46*cos(36)},{0.46*sin(36)}) {$\tfrac{2\pi}{5}$};
  % vertices and labels
  \filldraw[black] (z0) circle (0.02) node[right=2pt] {$z_0=1$};
  \filldraw[black] (z1) circle (0.02) node[above right=1pt] {$z_1=e^{2\pi i/5}$};
  \filldraw[black] (z2) circle (0.02) node[above left=1pt] {$z_2=e^{4\pi i/5}$};
  \filldraw[black] (z3) circle (0.02) node[below left=1pt] {$z_3=e^{6\pi i/5}$};
  \filldraw[black] (z4) circle (0.02) node[below right=1pt] {$z_4=e^{8\pi i/5}$};
  % title note (semi-transparent box) bottom-left, away from pentagon
  \node[fill=white,fill opacity=0.72,text opacity=1,draw=gray!40,rounded corners,
        align=center,font=\small] at (-0.92,-1.18)
        {$z^5=1$ 的五个根\\构成正五边形};
\end{tikzpicture}
\end{document}
